\section{A Worked Example: Permutation Circuits}\label{sec:permutations}

We begin with the library's smallest complete case study: circuits built
from a single two-wire gate, the swap $\sigma$, which present the symmetric
groups.  Everything in this section is verified; the statements quoted at
the end appear verbatim in \texttt{MainTheorems.agda}.  The point of the
walkthrough is to exhibit the pipeline --- gates, relations, cosets, tower,
semantics, uniqueness, completeness, presentation --- that every later case
study instantiates.

\subsection{Gates and circuits}

A gate set is a family $\aid{Gate} : \mathbb{N} \to \mathsf{Set}$ giving,
for each arity, the gates that act on that many adjacent wires.  For
permutations there is exactly one gate, on two wires:

\begin{agdacode}
data Gate : ℕ → Set where
  σ-gate : Gate 2
\end{agdacode}

The generic circuit layer (\S\ref{sec:circuits}) turns any such family into
wire-indexed \emph{generators} $\aid{Gen}\;n$ --- a gate placed at the bottom
of $n$ wires, or a generator shifted up one wire by the constructor
\aid{{\_}↥} --- and defines an $n$-wire circuit as a word of generators,
$\aid{Circuit}\;n = \aid{Word}\;(\aid{Gen}\;n)$.  Words (\S\ref{sec:words})
are the free monoid: a circuit is a generator $[\,g\,]$\aid{ʷ}, the empty
circuit \aid{ε}, or a composition $c \aid{•} c'$.  We write
$\sigma : \aid{Circuit}\;(2{+}n)$ for the one-letter circuit
$[\,\aid{gate₂}\;\aid{σ-gate}\,]$\aid{ʷ} and $c{\uparrow}$ for the circuit
$c$ shifted up one wire ($\aid{{\_}↑} = \aid{wmap}\;\aid{{\_}↥}$).

The reading is diagrammatic.  On three wires, the two generators and a
composite correspond to circuit diagrams as follows (wires drawn bottom-up,
composition left to right):

\begin{figure}[h]
\centering
\begin{minipage}[c]{0.2\linewidth}\centering
\begin{quantikz}[row sep=0.18cm, column sep=0.35cm]
  \lstick{2} & \qw       & \qw \\
  \lstick{1} & \swap{1}  & \qw \\
  \lstick{0} & \targX{}  & \qw
\end{quantikz}\\[2pt]
$\sigma$
\end{minipage}\hfill
\begin{minipage}[c]{0.2\linewidth}\centering
\begin{quantikz}[row sep=0.18cm, column sep=0.35cm]
  \lstick{2} & \swap{1}  & \qw \\
  \lstick{1} & \targX{}  & \qw \\
  \lstick{0} & \qw       & \qw
\end{quantikz}\\[2pt]
$\sigma{\uparrow}$
\end{minipage}\hfill
\begin{minipage}[c]{0.42\linewidth}\centering
\begin{quantikz}[row sep=0.18cm, column sep=0.35cm]
  \lstick{2} & \qw      & \swap{1} & \qw      & \qw \\
  \lstick{1} & \swap{1} & \targX{} & \swap{1} & \qw \\
  \lstick{0} & \targX{} & \qw      & \targX{} & \qw
\end{quantikz}\\[2pt]
$\sigma \mathbin{\bullet} \sigma{\uparrow} \mathbin{\bullet} \sigma$
\end{minipage}
\caption{Circuit code and diagrams for three-wire permutation circuits:
the Agda values \aid{σ}, \aid{σ ↑}, and \aid{σ • σ ↑ • σ} of type
\aid{Circuit 3}.}
\label{fig:sigma}
\end{figure}

\subsection{Relations: two axioms, structure for free}

The group-specific relations are an inductive family indexed by the wire
count --- note that each axiom is stated at the \emph{lowest} wire count at
which it makes sense:

\begin{agdacode}
data _SRel,_===_ : (n : ℕ) → WRel (Gen n) where
  order       : (₂₊ n) SRel, σ • σ === ε
  yang-baxter : (₃₊ n) SRel, σ • σ ↑ • σ === σ ↑ • σ • σ ↑
\end{agdacode}

These two axioms are the Coxeter presentation of $S_n$ in circuit clothing:
$\sigma^2 = \varepsilon$ and the braid (Yang--Baxter) relation.  What about
the relations every circuit theory needs --- that equations may be shifted
up wires, and that gates acting on disjoint wires commute?  Writing them
into every gate set would be repetitive and error-prone; the circuit layer
adds them generically.  \aid{Lift-Relation} (\S\ref{sec:circuits}) extends
any gate-specific family with three structural rules, yielding the full
relation set \aid{{\_}VRel,{\_}==={\_}}:

\begin{agdacode}
data _VRel,_===_ : (n : ℕ) → CRel n where
  srel  : n SRel, w === v → n VRel, w === v
  cong↑ : n VRel, w === v → (₁₊ n) VRel, w ↑ === v ↑
  comm₁ : (h : Gate 1) (g : Gen n) → (₁₊ n) VRel,
    [ g ↥ ]ʷ • [ gate₁ h ]ʷ === [ gate₁ h ]ʷ • [ g ↥ ]ʷ
  comm₂ : (h : Gate 2) (g : Gen n) → (₂₊ n) VRel,
    [ g ↥ ↥ ]ʷ • [ gate₂ h ]ʷ === [ gate₂ h ]ʷ • [ g ↥ ↥ ]ʷ
\end{agdacode}

The far-commutation instance of \aid{comm₂} and the braid relation are the
two diagram equalities of Figure~\ref{fig:relations}.  From here on,
$\approx$ denotes the monoid congruence generated by
$n\;\aid{VRel},{=}{=}{=}$ (\S\ref{sec:presentations}): the least relation
containing the axioms and closed under reflexivity, symmetry, transitivity,
congruence of \aid{•}, associativity, and the unit laws.  The word problem
for the presented monoid is precisely $\approx$.

\begin{figure}[h]
\centering
\begin{minipage}[c]{0.44\linewidth}\centering
\begin{quantikz}[row sep=0.18cm, column sep=0.3cm]
  & \swap{1} & \qw      & \qw \\
  & \targX{} & \qw      & \qw \\
  & \qw      & \swap{1} & \qw \\
  & \qw      & \targX{} & \qw
\end{quantikz}
$=$
\begin{quantikz}[row sep=0.18cm, column sep=0.3cm]
  & \qw      & \swap{1} & \qw \\
  & \qw      & \targX{} & \qw \\
  & \swap{1} & \qw      & \qw \\
  & \targX{} & \qw      & \qw
\end{quantikz}\\[2pt]
{\small \aid{comm₂}: disjoint gates commute}
\end{minipage}\hfill
\begin{minipage}[c]{0.52\linewidth}\centering
\begin{quantikz}[row sep=0.18cm, column sep=0.3cm]
  & \qw      & \swap{1} & \qw      & \qw \\
  & \swap{1} & \targX{} & \swap{1} & \qw \\
  & \targX{} & \qw      & \targX{} & \qw
\end{quantikz}
$=$
\begin{quantikz}[row sep=0.18cm, column sep=0.3cm]
  & \swap{1} & \qw      & \swap{1} & \qw \\
  & \targX{} & \swap{1} & \targX{} & \qw \\
  & \qw      & \targX{} & \qw      & \qw
\end{quantikz}\\[2pt]
{\small \aid{yang-baxter}}
\end{minipage}
\caption{The generic far-commutation rule and the gate-specific braid
relation as diagram equalities.}
\label{fig:relations}
\end{figure}

Because $\sigma$ is its own inverse and shifts of invertible generators are
invertible, the presented monoid is a group; the library witnesses this by
a \aid{Grouplike} proof (every generator has a left inverse), constructed by
induction on generators.

\subsection{Cosets and the staircase normal form}

The engine behind the normal form is coset enumeration along the subgroup
chain
\[
  S_1 \;<\; S_2 \;<\; \dots \;<\; S_n \;<\; S_{n+1},
\]
where $S_n$ sits inside $S_{n+1}$ as the subgroup fixing the bottom wire
--- syntactically, as the circuits of the form $c{\uparrow}$.  The
right cosets of $S_n$ in $S_{n+1}$ are indexed by where the bottom wire
goes: the identity coset, or one of $n$ cosets reached by a final descending
staircase of swaps.  In Agda, cosets and their \emph{section} back into
circuits are:

\begin{agdacode}
data C : ℕ → Set where          -- |C n| = n + 1
  ε   : C n
  σ•_ : C n → C (₁₊ n)

[_]ᶜ : C n → Circuit (₁₊ n)     -- the staircase section
[ ε ]ᶜ     = ε
[ σ• c ]ᶜ  = σ • ([ c ]ᶜ ↑)
\end{agdacode}

The whole level is driven by one small function, the \emph{coset table}: a
right action that pushes one generator past a coset, emitting a residual
circuit one level down and a successor coset:

\begin{agdacode}
ract : C (₁₊ n) → Gen (₂₊ n) → Circuit (₁₊ n) × C (₁₊ n)
ract ε         σ-gen  = ε , σ• ε
ract (σ• ε)    σ-gen  = ε , ε
ract (σ• σ• c) σ-gen  = σ , σ• σ• c
ract ε         (g ↥)  = [ g ]ʷ , ε
ract (σ• c)    (g ↥)  = proj₁ (ract c g) ↑ , σ• (proj₂ (ract c g))
\end{agdacode}

(The two remaining clauses handle vacuous low-arity cases.)  The library's
stateful word traversal \aid{{\_}ᵗ} (\S\ref{sec:words}) extends \aid{ract}
to whole circuits, threading the coset from left to right and concatenating
the residuals; the key soundness lemma, proved once per table, certifies
each transition against the congruence:
\[
  [\,c\,]^{\mathsf c} \aid{•} [\,b\,]\aid{ʷ}
  \;\approx\;
  b'{\uparrow} \aid{•} [\,c'\,]^{\mathsf c}
  \qquad\text{where } (b',c') = \aid{ract}\;c\;b .
\]
Running the extended action from the identity coset maps an $(n{+}2)$-wire
circuit to a pair of an $(n{+}1)$-wire circuit and a coset --- one level of
normalization.  \emph{Iterating} this is the coset tower
(\S\ref{sec:engine}): the Symmetric development instantiates the library's
\aid{CosetTower} with the family of presentations
$n \mapsto (n\;\aid{VRel},{=}{=}{=})$, the coset families $\aid{C}$, and one
\aid{Extension} record per level packaging \aid{ract}, the section, and the
five coset-table hypotheses.  The tower then folds a base normal form
upward:

\begin{agdacode}
NF : ℕ → Set
NF 0       = ⊤
NF (suc k) = T.tower-carrier ⊤ k    -- ⊤ × C 1 × C 2 × ⋯ × C k

nfp'-t : ∀ n → NormalForm (_VRel,_===_ n) (NF n)
nfp'-t 0       = base0'
nfp'-t (suc k) = T.nfp'-tower ext base1' k
\end{agdacode}

The carrier $\aid{NF}\;n$ has exactly $2 \cdot 3 \cdots n = n!$ elements:
a tuple of ``digits'', the $k$-th drawn from the $k{+}1$ cosets of
$S_k < S_{k+1}$.  This is the factorial number system, and it is no
accident --- \S\ref{sec:preliminaries} explains the general mixed-radix
pattern.  A \aid{NormalForm} witness (\S\ref{sec:nfproperty}) packages the
normal-form function together with its section
$\aid{inv-nf} : \aid{NF}\;n \to \aid{Circuit}\;n$ --- concretely, the
concatenation of the staircases $[\,c_k\,]^{\mathsf c}$, suitably shifted
--- and the proof $\aid{inv-nf} \circ \aid{nf} \approx \mathrm{id}$.
Figure~\ref{fig:staircase} shows a normal form's section: a ``staircase of
staircases'', exactly Selinger's normal form picture for
Clifford circuits restricted to swaps~\cite{selinger2015clifford}.

\begin{figure}[h]
\centering
\begin{quantikz}[row sep=0.16cm, column sep=0.3cm]
  \lstick{3} & \qw      & \qw      & \swap{1} & \qw      & \qw \\
  \lstick{2} & \qw      & \swap{1} & \targX{} & \swap{1} & \qw \\
  \lstick{1} & \swap{1} & \targX{} & \qw      & \targX{} & \qw \\
  \lstick{0} & \targX{} & \qw      & \qw      & \qw      & \qw
\end{quantikz}
\caption{The section of the normal form
$(\mathsf{tt},\, \sigma{\bullet}\varepsilon,\,
\sigma{\bullet}\sigma{\bullet}\varepsilon,\,
\sigma{\bullet}\varepsilon) \in \aid{NF}\;4$: reading right to left, one
descending staircase per tower level.}
\label{fig:staircase}
\end{figure}

\subsection{Two semantics, uniqueness, and completeness}

So far everything was syntax.  The development interprets circuits twice:

\begin{itemize}
\item a \emph{loose} semantics
  $\aid{⟦\_⟧} : \aid{Circuit}\;n \to (\mathsf{Fin}\;n \to \mathsf{Fin}\;n)$
  into endofunctions of $\mathsf{Fin}\;n$ under composition ($\sigma$ swaps
  the bottom two elements), a monoid that is \emph{not} a group; and
\item a \emph{tight} semantics into
  $\aid{Permutation′-group}\;n$, the group of permutations of
  $\mathsf{Fin}\;n$ --- the standard library's \aid{Permutation′} bundled
  as a \aid{Group} by our \texttt{ForStdlib} layer
  (\S\ref{sec:forstdlib}).
\end{itemize}

\emph{Soundness} --- $w \approx v$ implies
$\aid{⟦}w\aid{⟧} \doteq \aid{⟦}v\aid{⟧}$ --- is, as always, a small case
analysis over the axioms: the two Coxeter relations hold for the swap
function, and the structural rules hold because shifted circuits act only on
shifted points.  The heart of the matter is the \emph{unique normal form}
property (\S\ref{sec:nfproperty}): normal forms with equal denotations are
equal,
\[
  \aid{⟦}\,\aid{inv-nf}\;u\,\aid{⟧} \doteq \aid{⟦}\,\aid{inv-nf}\;v\,\aid{⟧}
  \;\longrightarrow\; u \equiv v ,
\]
proved by induction on the tower: two staircase sections that agree as
functions must peel off equal top digits (they send the top point to the
same place), and the rest agree by induction.  Soundness and uniqueness are
exactly the premises of the library's generic lemma
\aid{by-normalization}, which concludes \emph{completeness} --- injectivity
of the semantics on congruence classes.  The three theorems, as restated in
\texttt{MainTheorems.agda}:

\begin{agdacode}
symmetric-soundness :
  ∀ n → Congruent (PB._≈_ (n VRel,_===_))
          (Setoid._≈_ (SymLoose.Endo-setoid n)) (SymLoose.⟦_⟧ {n})

symmetric-completeness :
  ∀ n → Injective (PB._≈_ (n VRel,_===_))
          (Setoid._≈_ (SymLoose.Endo-setoid n)) (SymLoose.⟦_⟧ {n})

symmetric-unique-nf :
  ∀ n → NFBase.UniqueNormalForm (n VRel,_===_) (SymNF.NF n)
          (Group.setoid (SymTight.Permutation′-group n))
          (SymTight.⟦_⟧ {n}) (SymNF.nfp'-t n)
\end{agdacode}

Completeness for the \emph{loose} semantics is the stronger, perhaps
surprising fact: even read as mere endofunctions, circuits equal as
functions are provably equal --- the relation set does not need
invertibility of the model.  For the tight semantics, soundness,
completeness, and surjectivity (every permutation is a product of adjacent
transpositions --- realized by \aid{inv-nf} composed with a semantic normal
form) assemble into the presentation record (\S\ref{sec:presentations}):

\begin{agdacode}
symmetric-presentation :
  ∀ n → (n VRel,_===_) IsPresentationOf (SymTight.Permutation′-group n)
\end{agdacode}

read: the free group on swap circuits modulo \{\aid{order},
\aid{yang-baxter}\} + structural rules \emph{is} the group of permutations
of $\mathsf{Fin}\;n$, via a group isomorphism whose underlying map is the
semantics.  Everything later in the paper is this pipeline again, with
richer gates, longer towers, and amalgamations where the chain branches.
